\chapter{El Balance}
\section{El balance}

\section{El activo}
\paragraph{Activo: deficinión y requisitos reconocimiento.}
Derechos y otros recursos controlados económicamente por la empresa, resultates de sucesos pasados, de los que se espera que la empresa obtenga beneficios o rendimientos económicos en el futuro.

Los tres atributos para que una partida adquiera la condición de activo:
\begin{itemize}
    \item Que la empresa controle o sea beneficiaria de los rendimientos esperados de los recursos
    \item Que el recurso se haya obtenido en virtud de alguna transferencia o circunstancia ya acontecida
    \item Que sea alta la probabilidad de obtener beneficios en un futuro
\end{itemize}

\subsection{El activo no corriente.}
\paragraph{El activo no corriente y su valoración.}
En general, el convenio para el tratamiento contable del inmovilizado material o inmateria, es el siguiente:
\begin{itemize}
    \item valoración a la compra, que se efectúa atendiedo al criterio del precio de adquisición
    \begin{enumerate}
        \item Se valora por el precio facturado por el vendedor más los astos de registro y otros gastos adicionales.
        \item Se valora por el coste de producción y en éste se incluyen los costes en que hubiera incurrido para su desarrollo. Se entiende que un elemento está creado o desarrollado cuando sea viable técnicamente y rentable económicamente.
    \end{enumerate}
    \item valoración al cierre, que exige reconocer una pérdida por deterioro del activo, siempre que el valor recuperable del activo sea menor que el valor neto contable
    \item valor a la venta o retiro, que exige reconocer la pérdida o ganancia en la cuenta de resultados
\end{itemize}

\paragraph{El activo no corriente. El inmovilizado inmaterial.}
Los activos intangibles son aquellos recursos obtenidos por una empresa que no son de naturaleza material. Los activos intangibles proveen derecho o privilegio a la empresa que los posee y son habitualmente adquiridos para uso operacional. Tipos de inmovilizado intangible:
\begin{itemize}
    \item Gastos de investigación y desarrollo
    \item Concesiones administrativas
    \item Propiedad industrial
    \item Fondo de comercio
    \item Derechos de traspaso
    \item Aplicaciones informáticas
\end{itemize}

Los activos intangibles deber ser identificables, es decir:
\begin{itemize}
    \item Ser susceptible de ser vendidos, arrendados, intercambiados.
    \item Surgen de derechos legales o contractuales
\end{itemize}

No son activos intangibles:
\begin{itemize}
    \item Marcas
    \item Listas de clientes
    \item Otras partidas generadas internamente
\end{itemize}

Con el paso del tiempo estos elementos ven disminuidas sus posibilidades de generar ingresos en el futuro. De ahí que los activos intangibles se amortizan durante la vida útil estimada, para que se pueda reconocer su contribución y tendrán una vida útil indefinida cuando se espere que el activo continúe generando flujos de efectivo para la empresa.
Se analizan su deterioro cuando existan indicios de eventual deterioro y al menos anualmente.
Se debe optar por amortizarlos en el período más corto posible, lo más rápidamente posible. Nuestra legislación establece ciertaws limitaciones en cuanto a plazos de amortización de estos bienes.

\subsection{El activo corriente.}
\subsubsection{Las cuentas de clientes. Préstamos y partidas a cobrar por operaciones comerciales.}
\paragraph{Valoración inicial.} Al valor razonable de la contraprestación entregada más los gastos de transacción que les sean directamente atribuibles.

\paragraph{Valoración psoterior.} Las partidas por cobrar a corto plazo y en las que no exista una tasa de interés contractual, se valoran según importe de factura. 

El resto de activos se valorarán por su coste amortizado y los intereses devengados se contabilizarán en la cuenta de pérdidas y ganancias.

Al cierre del ejercicio deberán hacerse las correcciones valorativas necesarias siempre que exista evidencia objetiva de su deterioro y que se producirá una reducción del crédito.

\subsubsection{Las cuentas de clientes. El descuento de efectos}
Alternativas de cobro de los efectos comerciales:
\begin{itemize}
    \item Esperar al vencimiento para cobrar el importe.
    \item Descontar el efecto en una entidad financiera. El banco pone a disposición de la empresa el importe del efecto antes de su vencimiento previa deducción de unos gastos. La empresa mejora su liquidez sin esperar al vencimiento de sus derechos de cobro. La operación genera una deuda subsidiaria con el banco que adelanta el importe.
    \item Encargar la gestión de cobro a una entidad. La entidad bancaria actúa como agente cobrado y por ello percibe una comisión.
    \item Traspasar los derechos de cobro a otra entidad o tercero (factoring sin recurso). La empresa vende, antes de su vencimiento, el effecto a una sociedad de factoring, que se encarga de su cobro y que puede asumir el riesgo de impago o cambio de una comisión mayor.
    \item Confirming. Una fórmula financiera a través de la cual una empresa contrata a una sociedad financiera para que sea esta quien se encargue de la gestión de sus pagos a proveedores.
\end{itemize}

\section{El pasivo y el patrimonio neto}
\subsection{El patrimonio neto}
\paragraph{Fondos propios}
\begin{itemize}
    \item Capital
    \item Prima de emisión
    \item Reservas
    \item Resultado de ejercicios anteriores
    \item (Dividendo a cuenta)
\end{itemize}
\paragraph{Ajustes por cambios de valor}
\paragraph{Subvenciones, donaciones y legados recibidos}
\paragraph{Capital social}
Garantiza la solvencia o responsabilidad de la empresa frente a terceros y acreedores. Los socios no responden personalmente de las deudas sociales.
\[ Capital\ social = n^\circ\ acciones \times valor\ nominal \]
Las acciones representan la parte (alícuota) del capital social que representa la propiedad de los socios en la empresa. 
Las acciones otorgan a los propietarios derechos económicos, políticos y de información.
\begin{itemize}
    \item Derechos políticos, derecho a participar en la toma de decisiones en la junta general de accionistas.
    \item Derechos económicos, derecho a percibir ``dividendos'', es decir, a recibir parte del beneficio generado por la empresa, derecho preferente de suscripción de nuevas acciones ante una ampliación de capital.
    \item Derechos de información, derecho a recibir información periódica sobre la marcha de la sociedad.
\end{itemize}

\paragraph{Resultado del ejercicio}
Informa de la marcha económica de la empresa, beneficio o pérdida, y se aplica a:
\begin{itemize}
    \item Dotar reservas obligatorias (legal, estatuaria, fondo de comercio...) y voluntarias.
    \item Compensación de pérdidas de ejercicios anteriores
    \item Distribución de dividendos, reparto del beneficio a los socios.
\end{itemize}

\paragraph{Reservas}
Son cuentas de Patrimonio Neto que representan recursos propios acumulados y según sea posible su libre disposición o esté restringido su uso:
\begin{itemize}
    \item Reserva legal: Obligatoria, un $10 \%$ del beneficio se destina a la reserva legal hasta que ésta alcance al menos el $20 \%$ del capital social.
    \item Reservas especiales: Establecidas por disposición obligatoria. Dotación con cargo a ``Pérdidas y Ganancias''
    \item Reservas voluntarias y estatutarias: Su dotación es libre, cubiertos los requisitos legales, se puede dotar estas reservas en lugar de distribuir el resultado.
    \item Reserva para acciones propias: Si existe autocartera y en determinadas condiciones.
    \item Reservas especiales: por fondo de comercio, I+D \dots
\end{itemize}


\subsection{El pasivo}
\subsubsection{Pasivos corrientes}
Los pasivos corrientes (circulantes o exigibles a C/P) financian en parte el activo corriente de la empresa.
Su plazo de reembolso es menor que un año o el ciclo de explotación.
\paragraph{Acreedores y deudores por operaciones comerciales.}
\begin{itemize}
    \item Proveedores
    \item Acreedores varios
    \item Personal
    \item Administración pública
    \item Ingresos anticipados (ajustes por periodificación)
\end{itemize}

\paragraph{Cuentas financieras}
\begin{itemize}
    \item Emprésitos, deudas con características especiales a C/P
    \item Deudas a C/P con partes vinculadas
    \item Deudas a C/P por préstamos recibidos y otros conceptos
    \item Fianzas y depósitos recibidos a C/P y ajustes por periodificación (gasto de intereses)
\end{itemize}

\subsubsection{Préstamos y partidas a pagar por operaciones comerciales}
Representan una deuda u obligación por el aplazamiento en el pago en la adquisición de bienes y servicios utilizados en el proceso productivo o por la compra de activo no corriente.
\begin{itemize}
    \item Proveedores, deudas asociadas a la actividad principal de la empresa
    \item Acreedores, deudas no relacionadas con la actividad principal de la empresa
    \item Otras: proveedores de inmovilizado a c/p, deudas con el personal y deudas con las administraciones públicas
\end{itemize}

\paragraph{Valoración inicial.}
AL valor razonable de la contraprestación entregada más los gastos de transacción que les sean directamente atribuibles.

\paragraph{Valoración posterior.}
Las partidas a pagar a corto plazo y en las que no exista una tasa de interés contractual, se valoran según importe de factura. El resto de pasivos se valorarán por su coste amortizado y caso de que haya intereses asociados a la operación se contabilizaran siguiendo el método del interés efectivo.

\subsubsection{Provisiones y Contingencias}
Una contingencia es aquella obligación que surgida de acontecimientos pasados, no se reconocen contablemente en el acontecimineto presente porque:
\begin{itemize}
    \item No es probable que la entidad tenga que satisfacerla, desprendiéndose de recursos económicos que incorporen beneficios.
    \item El importe de la obligación no puede ser medido con la suficiente fiabilidad.
    \item Solo se confirmará si acontecen o no sucesos futuros inciertos que no están enteramente bajo el control de la empresa.
\end{itemize}

Los pasivos contingentes solo figurarán como dato informativo en la Memoria.

Una provisión se debe reconocer únicamente cuando se cumplen las siguientes condiciones:
\begin{enumerate}
    \item La entidad tiene una obligación presente surgida a raíz de sucesos pasados
    \item Es probable que para cancelarla la empresa deba desprenderse de recursos económicos y
    \item Puede efectuarse una estimación fiable de la deuda correspondiente
\end{enumerate}

Las provisiones son pasivos en el que existe incertidumbre acerca de su cuantía o vencimiento

\subsubsection{Dotación de la provisión}
Por el desembolso a efectuar en el futuro por la obligación contraída (tasada al valor actual mediente la aplicación de un tipo de descuento)
\begin{itemize}
    \item Impuestos
    \item Desmantelamiento, retiro o rehabilitación del inmovilizado
    \item Actuaciones medioambientales
    \item Reestructuraciones
    \item Otras responsabilidades como litigios, indemnizaciones\dots
\end{itemize}

Una vez reconocida la provisión, deber ser objeto de revisión al cierre. Los ajustes se contabilizarán como un gasto financiero según se vayan devengando.

\subsubsection{Préstamos}
\paragraph{Tipos de préstamos.}
\begin{itemize}
    \item Préstamos con entidades de crédito
    \item Préstamos con otras empresas (Ej.: empresas del grupo)
    \item Préstamos para la adquisición de activos NC (proveedores inmovilizado).
    \item Devolución del préstamo: según sistema americano, francés, italiano, etc.
\end{itemize}

\paragraph{Valoración inicial.}
Los préstamos se reconocerán inicialmente por su valor razonable, (precio de la transacción), es decir, el valor razonable de la contraprestación recibida ajustado por los costes de la transacción directamente atribuibles.

\paragraph{Valor posterior.}
Con posterioridad, los préstamos quedarán valorados a su coste amortizado: importe inicial, más los intereses devengados menos los pagos realizados (por intereses o devolución de parte del principal).

\section{El equilibrio financiero}

